\clearpage
\color{red}
\textbf{TODO:}

\subsection{Lag}



\subsection{Rolling}


\subsection{Feature Sicherung}


\subsection{Robustheit}
\itemize{
    \item Robustheit gegenüber unvollständigen Daten
    \item Robustheit gegenüber fehlerhaften Daten
    \item Robustheit gegenüber unregelmäßigen Daten
}

\begin{itemize}
    \item Drop von Datenpunkten bei fehleraften werten
\end{itemize}

\paragraph{2026-05-26 --- Modul-Reorganisation in \texttt{forecast/} und \texttt{classifier/}}
Die beiden Modellfamilien (binärer Driving-Classifier vs.\ 1-7-Tage-Forecaster)
wurden in eigene Unterordner unter \texttt{code/model\_scripts/} verschoben.
Geteilte Helpers (\texttt{base.py}, \texttt{data\_adapters.py}) bleiben oben.
Begründung: klare Trennung der Forschungsschienen --- die Cross-Validation-
Methodik (Datenqualität, binär) hängt nicht mit der Anwendungs-Schiene
(Tages-Forecast, multi-vehicle, Quantile) zusammen.

\paragraph{2026-05-26 --- Ablations-Registry und \texttt{LGBMForecaster}}
Eingeführt: \texttt{forecaster\_registry.py} mit Lazy-Import-Lookup
\texttt{name $\to$ Klasse}. Erste Schwesterimplementierung
\texttt{LGBMForecaster} mit nativer Quantil-Regression (drei separate
LGBMRegressoren mit \texttt{objective="quantile"}, \texttt{alpha=0.1/0.5/0.9})
statt RF-Tree-Variance. Damit ist die Ablationsstudie
(LogReg $\to$ RF $\to$ LightGBM $\to$ LSTM $\to$ TFT) erweiterbar, ohne
Backend, Trainingsskripte oder Frontend zu ändern. Backend-Schnittstelle
auf \texttt{predict\_proba\_day} + \texttt{predict\_quantiles} abstrahiert
--- keine Direkt-Inspektion von \texttt{reg.estimators\_} mehr.

\paragraph{2026-05-26 --- Multi-Vehicle-Forecaster}
\texttt{RandomForestForecaster.fit} akzeptiert jetzt einen
\texttt{vehicle\_id}-mehrwertigen Eingabe-Frame. Trainingszeilen kommen aus
allen Fahrzeugen (Lag-Fenster überschreiten nie Fahrzeug-Grenzen). Für
\texttt{predict()} wird das aktivste Fahrzeug mit jüngstem
\texttt{timestamp.max()} als Default ausgewählt; dessen Daily-Slice wird
in \texttt{self.daily} persistiert. Das 7$\times$24-Stundenprofil mittelt
mit gleichem Gewicht über alle Fahrzeuge, damit langlebigere Zeitreihen
nicht dominieren. Erlaubt einen einzigen Forecaster pro Multi-Vehicle-Quelle
(\texttt{emobpy}, \texttt{ved}).

\paragraph{2026-05-26 --- Run-Verzeichnis-Schema}
Jede Forecast-Erzeugung schreibt in
\texttt{predictions/<modell>/<datensatz>\_<N>d\_T<DD-MM-YYYY>/}. Ein Run
ist damit ein vollständiges Verzeichnis-Objekt (CSV, Plots, Notizen) und
kann über das Frontend in den Store importiert werden. Hinweis: Doppelpunkte
sind in Windows-Pfaden verboten, deshalb \texttt{T<DD-MM-YYYY>} mit
Bindestrich statt \texttt{TDD:MM:YYYY}.

\paragraph{2026-05-26 --- Forecaster-Familie pro Datenquelle}
Trainingsskripte
\texttt{train\_<source>\_forecaster.py} für alle vier Quellen ergänzt.
\texttt{real\_world\_ev} ersetzt das frühere \texttt{car\_full\_forecaster}
unter konsistentem Namensschema \texttt{<source>\_forecaster\_<algo>.joblib}.
Das alte joblib bleibt für Backward-Compat erhalten; ein Shim in
\texttt{model\_scripts/csv\_forecaster.py} re-exportiert die Klasse unter
dem alten Modulpfad, damit \texttt{joblib.load} weiterhin funktioniert.

\color{black}

\clearpage